\section{Introduction}

Modern processors often stall because of dependencies between
instructions. Value prediction tries to reduce
this delay by guessing an instruction's result before it is
available, so that dependent operations can start earlier. A wrong
guess must be detected and handled without breaking correctness.

This project builds on
``Architecting Value Prediction around In-Order Execution''
\cite{ravenel2025architecting}. The paper studies value prediction in
in-order processors using Last Value Prediction (LVP), ASTRA-based
load-base prediction, and V-ASTRA. V-ASTRA combines a predicted load
base with a dedicated value cache, so load values are available
earlier. The paper also introduces META, which coordinates LVP and
V-ASTRA. It further describes address compression, store tracking,
scoreboard replay, and a speculative L0 minicache.

My objective was to implement and evaluate similar mechanisms in
gem5 using MinorCPU. MinorCPU differs from the processor in the
original study, so I had to change how it handles loads, their
dependent instructions, and recovery from wrong predictions.

I implemented an LVP predictor and an ASTRA-based predictor for load
base addresses. V-ASTRA adds the instruction's offset to the predicted
base and reads the load value from a small value cache. When a set of
safety checks passes, the load can skip its normal memory request.

On top of these, I tried two ways of letting dependent instructions start
earlier. Design~A only releases the scoreboard dependency, while
Design~B also forwards the validated value. I also adapted META to
choose at most one of LVP or V-ASTRA for each dependent load. Finally,
an experimental option lets loads served by LVP feed their base
addresses to ASTRA, both for training and for V-ASTRA candidacy.

I evaluated the implementation using directed RISC-V regression
tests and deterministic microbenchmarks. The tests check that predictions,
bypasses, and recovery behave correctly. The performance experiments
compare no prediction (NP), LVP, V-ASTRA, and META on workloads with
different address, value, and dependency patterns, including the
experimental provisional-base modes.
